\documentclass[11pt,a4paper]{article}

% ---------- Fonts & typography ----------
\usepackage[T1]{fontenc}
\usepackage[utf8]{inputenc}
\usepackage{newpxtext,newpxmath}
\usepackage[protrusion=true]{microtype}
\linespread{1.02}

% ---------- Page layout ----------
\usepackage[a4paper,left=2.1cm,right=1.5cm,top=1.7cm,bottom=1.8cm]{geometry}
\setlength{\parindent}{0pt}
\setlength{\parskip}{0.45em}

% ---------- Colours ----------
\usepackage{xcolor}
\definecolor{accent}{RGB}{28,58,104}
\definecolor{subtle}{RGB}{110,110,110}

% ---------- Lists ----------
\usepackage{enumitem}
\setlist[itemize]{leftmargin=1.2em,topsep=0.15em,itemsep=0.2em,parsep=0pt}

% ---------- Section headings ----------
\usepackage{titlesec}
\titleformat{\section}
  {\normalsize\bfseries\color{accent}}
  {}{0pt}{}
\titlespacing*{\section}{0pt}{0.7em}{0.15em}

\pagestyle{empty}

% ---------- Links ----------
\usepackage[hidelinks]{hyperref}
\hypersetup{
  pdftitle={Research Plan},
  pdfauthor={Avi Kadria}
}

\begin{document}

% ---------- Title block ----------
{\large\bfseries\color{accent}Research Plan}\hfill Avi Kadria\\[0.1em]
\null\hfill{\small\color{subtle}\href{mailto:avi.kadria3@gmail.com}{avi.kadria3@gmail.com}}
\par\vspace{0.3em}

My research combines graph algorithms and fine-grained complexity, with a focus on distance problems: shortest paths, distance oracles, routing schemes, girth, and diameter. These problems trade off running time, space, and approximation quality. My goal is to determine the optimal trade-offs, both by designing faster algorithms and by proving lower bounds that show when no further improvement is possible. My work has appeared at SODA, ICALP, ESA, ISAAC (Best Paper Award), and DISC. As a postdoc, I plan to pursue two connected directions.

\section{1. Faster Algorithms for Distance Problems}

My recent results improve several state-of-the-art bounds. With Roditty, I gave the first improvement on the construction time of the Thorup--Zwick distance oracles for graphs that are not sparse, including the first subquadratic construction for $2<k<6$, which resolves an open problem of Wulff-Nilsen [ICALP'26]. We also gave the first improvement for single-source shortest paths (SSSP) in weighted undirected graphs [SOFSEM'27] since the breakthrough of Duan, Mao, Shu, and Yin [FOCS'23]. For $n$-Pairs Shortest Paths, we introduced the \emph{heavy-edge} technique and obtained a $1.622k$-approximation in weighted graphs, answering an open question of Chechik, Hoch, and Lifshitz [ESA'26]. I plan to continue in four directions:
\begin{itemize}
  \item \textbf{Distance oracles.} How fast can $(2k-1)$-stretch oracles be constructed? Wulff-Nilsen asked for near-linear constructions; we extended the range of densities where they are possible. 
  I aim to close the remaining gap by refining the sampling and exploration techniques underlying our parameterized oracle hierarchy.
  \item \textbf{Shortest paths.} Our SSSP algorithm shows that the $O(m\sqrt{\log n\log\log n})$ bound is not the limit. I aim to determine the true complexity of SSSP in the comparison-addition model, and whether our improvement extends to directed graphs, where Duan et al.\ [ICALP'26] recently matched the undirected bound.
  \item \textbf{Derandomization.} Most fast distance algorithms rely on random sampling. My deterministic girth algorithm (arXiv'26) matches the randomized bounds through a new way to truncate the Thorup--Zwick hierarchy. I plan to apply this idea to distance oracles, spanners, and $n$-Pairs Shortest Paths.
  \item \textbf{Geometry and graph structure.} I plan to use geometric techniques to design faster distance algorithms for structured graph classes. In ongoing work on weighted directed planar graphs, we reduce $(1+\varepsilon)$-approximate diameter to orthogonal range queries and obtain near-linear time.
\end{itemize}

\section{2. Lower Bounds and Structural Barriers}

For sparse graphs, the classical space lower bound for distance oracles gives nothing beyond $\Omega(m)$. With Roditty, I proved that every pivot-based oracle with stretch $k$ requires $\tilde\Omega(m^{1+2/(k+1)})$ space on sparse graphs. This bound is tight, and the model captures all state-of-the-art oracles with $\tilde O(1)$ query time. We also proved new fine-grained lower bounds for girth approximation [ESA'26]. I plan to:
\begin{itemize}
  \item \textbf{Go beyond restricted models.} Extend space lower bounds from pivot-based oracles toward general oracles for sparse graphs, a central open question in the area.
  \item \textbf{Prove conditional lower bounds.} Use 3SUM, APSP, and related hypotheses to prove lower bounds for girth, shortest cycles, and $n$-Pairs Shortest Paths, and determine how the hardness depends on graph density and approximation factor.
\end{itemize}

\section{3. Long-Term Vision}

My long-term goal is to close the gaps between upper and lower bounds for fundamental graph problems. Algorithms and lower bounds inform each other: a faster algorithm rules out hardness arguments, and a tight lower bound shows which techniques cannot succeed. I aim to develop reusable tools, like the heavy-edge technique and the hierarchy of parameterized oracles, that explain why particular polynomial bounds arise and when they can be beaten.

\end{document}
